% GENERATED FILE. Edit canonical lessons, then run npm run generate:books.
% canonical-source-hash: fnv1a64:30bed3161b49856a
% canonical-lessons: HI-C385-chitrana, HI-C385-mitna, HI-C385-latarna, HI-C385-dhalna, HI-C385-adhura

\chapter{If Things Go Wrong}
\label{ch:hi-if-things-go-wrong}
\hlchaptermodality{385}

\begin{chapteropening}
\textbf{By the end of this chapter:} \emph{I can warn what will go wrong if something is done or left undone.}

Complete the last lesson of chapter 385: I can warn what will go wrong if something is done or left undone.
\end{chapteropening}

\section[\texorpdfstring{\hi{छितराना} (chitrānā)}{छितराना (chitrānā)}]{\hi{छितराना} (chitrānā) — to scatter, to strew}

\label{lesson:HI-C385-chitrana}



\begin{quote}
Before the new one: say the Hindi for to soften, then the Hindi for moderate.
\end{quote}

\subsection*{You'll want to know: \hi{छितराना}}
\textbf{\hi{छितराना}} — \emph{chitrānā} — \textquotedblleft{}to scatter, to strew\textquotedblright{}.

\begin{quote}
\hi{अगर} \hi{तुम} \hi{आँगन} \hi{में} \hi{चावल} \hi{छितराओगे} \hi{तो} \hi{चिड़ियाँ} \hi{आएँगी।}
\end{quote}

\begin{quote}
\emph{agar tum āṁgan meṁ cāval chitrāoge to ciṛiyā̃ āeṅgī.}
\end{quote}

\begin{quote}
If you scatter rice in the courtyard, the birds will come.
\end{quote}

Past, with \hi{ने} (he or she …): \hi{उसने} \hi{बीज} \hi{छितराए}, \emph{usne bīj chitrāe}.

\begin{cousinweb}
\textbf{\hi{छितराना}} (\emph{chitrānā}) is inherited from \emph{*chitr}, \textquotedblleft{}scatter, sprinkle\textquotedblright{}, a form CDIAL reconstructs from its descendants (CDIAL 5040).
\end{cousinweb}

\begin{grammarlens}[title={What you've built}]
One more word for saying what happens if.
\end{grammarlens}

\subsection*{Your turn}
\begin{itemize}
\raggedright
  \item \emph{Say it:} \emph{chitrānā}
  \item \emph{Say it:} \emph{chitrānā}, once more
  \item \emph{Say it:} \emph{agar tum āṁgan meṁ cāval chitrāoge to ciṛiyā̃ āeṅgī.} — If you scatter rice in the courtyard, the birds will come.
  \item \emph{Recall:} say the Hindi for to soften, then the Hindi for moderate, then say \emph{chitrānā} again
\end{itemize}

\begin{tcolorbox}[breakable,colback=teal!4,colframe=teal!35!black,title={Before you move on}]
What does \hi{छितराना} mean? (\textquotedblleft{}To scatter, to strew\textquotedblright{}.) Say it once more.
\end{tcolorbox}
\section[\texorpdfstring{\hi{मिटना} (miṭnā)}{मिटना (miṭnā)}]{\hi{मिटना} (miṭnā) — to be wiped out, to fade away}

\label{lesson:HI-C385-mitna}



\begin{quote}
Before the new one: say the Hindi for moderate, then the Hindi for to strew.
\end{quote}

\subsection*{You'll want to know: \hi{मिटना}}
\textbf{\hi{मिटना}} — \emph{miṭnā} — \textquotedblleft{}to be wiped out, to fade away\textquotedblright{}.

\begin{quote}
\hi{अगर} \hi{बारिश} \hi{होगी} \hi{तो} \hi{दीवार} \hi{का} \hi{रंग} \hi{मिट} \hi{जाएगा।}
\end{quote}

\begin{quote}
\emph{agar bāriś hogī to dīvār kā raṅg miṭ jāegā.}
\end{quote}

\begin{quote}
If it rains, the colour on the wall will fade.
\end{quote}

Past (it …): \hi{वह} \hi{नाम} \hi{मिट} \hi{गया}, \emph{vah nām miṭ gayā}.

\begin{cousinweb}
\textbf{\hi{मिटना}} (\emph{miṭnā}) is a tadbhava, inherited from Sanskrit \emph{mṛṣṭa}, \textquotedblleft{}rubbed, wiped\textquotedblright{} (CDIAL 10299).
\end{cousinweb}

\begin{grammarlens}[title={What you've built}]
One more word for saying what happens if.
\end{grammarlens}

\subsection*{Your turn}
\begin{itemize}
\raggedright
  \item \emph{Say it:} \emph{miṭnā}
  \item \emph{Say it:} \emph{miṭnā}, once more
  \item \emph{Say it:} \emph{agar bāriś hogī to dīvār kā raṅg miṭ jāegā.} — If it rains, the colour on the wall will fade.
  \item \emph{Recall:} say the Hindi for moderate, then the Hindi for to strew, then say \emph{miṭnā} again
\end{itemize}

\begin{tcolorbox}[breakable,colback=teal!4,colframe=teal!35!black,title={Before you move on}]
What does \hi{मिटना} mean? (\textquotedblleft{}To be wiped out, to fade away\textquotedblright{}.) Say it once more.
\end{tcolorbox}
\section[\texorpdfstring{\hi{लताड़ना} (latāṛnā)}{लताड़ना (latāṛnā)}]{\hi{लताड़ना} (latāṛnā) — to tell off, to scold harshly}

\label{lesson:HI-C385-latarna}



\begin{quote}
Before the new one: say the Hindi for to strew, then the Hindi for to be wiped out.
\end{quote}

\subsection*{You'll want to know: \hi{लताड़ना}}
\textbf{\hi{लताड़ना}} — \emph{latāṛnā} — \textquotedblleft{}to tell off, to scold harshly\textquotedblright{}.

\begin{quote}
\hi{अगर} \hi{तुमने} \hi{फिर} \hi{कमरा} \hi{गंदा} \hi{किया} \hi{तो} \hi{माँ} \hi{तुम्हें} \hi{लताड़ेगी।}
\end{quote}

\begin{quote}
\emph{agar tumne phir kamrā gandā kiyā to māṁ tumhẽ latāṛegī.}
\end{quote}

\begin{quote}
If you make the room dirty again, your mother will tell you off.
\end{quote}

Past, with \hi{ने} (he or she …): \hi{उसने} \hi{मुझे} \hi{लताड़ा}, \emph{usne mujhe latāṛā}.

\begin{cousinweb}
\textbf{\hi{लताड़ना}} (\emph{latāṛnā}) is inherited from \emph{*lattā}, \textquotedblleft{}a kick\textquotedblright{}, a form CDIAL reconstructs from its descendants (CDIAL 10931). CDIAL glosses the Hindi verb \textquotedblleft{}to spurn\textquotedblright{}; telling someone off, the everyday sense, is not among its glosses.
\end{cousinweb}

\begin{grammarlens}[title={What you've built}]
One more word for saying what happens if.
\end{grammarlens}

\subsection*{Your turn}
\begin{itemize}
\raggedright
  \item \emph{Say it:} \emph{latāṛnā}
  \item \emph{Say it:} \emph{latāṛnā}, once more
  \item \emph{Say it:} \emph{agar tumne phir kamrā gandā kiyā to māṁ tumhẽ latāṛegī.} — If you make the room dirty again, your mother will tell you off.
  \item \emph{Recall:} say the Hindi for to strew, then the Hindi for to be wiped out, then say \emph{latāṛnā} again
\end{itemize}

\begin{tcolorbox}[breakable,colback=teal!4,colframe=teal!35!black,title={Before you move on}]
What does \hi{लताड़ना} mean? (\textquotedblleft{}To tell off, to scold harshly\textquotedblright{}.) Say it once more.
\end{tcolorbox}
\section[\texorpdfstring{\hi{ढलना} (ḍhalnā)}{ढलना (ḍhalnā)}]{\hi{ढलना} (ḍhalnā) — to set (of the sun), to wane}

\label{lesson:HI-C385-dhalna}



\begin{quote}
Before the new one: say the Hindi for to be wiped out, then the Hindi for to tell off.
\end{quote}

\subsection*{You'll want to know: \hi{ढलना}}
\textbf{\hi{ढलना}} — \emph{ḍhalnā} — \textquotedblleft{}to set (of the sun), to wane\textquotedblright{}.

\begin{quote}
\hi{अगर} \hi{सूरज} \hi{ढल} \hi{जाएगा} \hi{तो} \hi{रास्ता} \hi{नहीं} \hi{दिखेगा।}
\end{quote}

\begin{quote}
\emph{agar sūraj ḍhal jāegā to rāstā nahīṃ dikhegā.}
\end{quote}

\begin{quote}
If the sun sets, we won't be able to see the path.
\end{quote}

Past (it …): \hi{वह} \hi{सूरज} \hi{ढल} \hi{गया}, \emph{vah sūraj ḍhal gayā}.

\begin{cousinweb}
\textbf{\hi{ढलना}} (\emph{ḍhalnā}) is inherited from \emph{*ḍhalati}, \textquotedblleft{}bends over, falls\textquotedblright{}, a form CDIAL reconstructs from its descendants (CDIAL 5581).
\end{cousinweb}

\begin{grammarlens}[title={What you've built}]
One more word for saying what happens if.
\end{grammarlens}

\subsection*{Your turn}
\begin{itemize}
\raggedright
  \item \emph{Say it:} \emph{ḍhalnā}
  \item \emph{Say it:} \emph{ḍhalnā}, once more
  \item \emph{Say it:} \emph{agar sūraj ḍhal jāegā to rāstā nahīṃ dikhegā.} — If the sun sets, we won't be able to see the path.
  \item \emph{Recall:} say the Hindi for to be wiped out, then the Hindi for to tell off, then say \emph{ḍhalnā} again
\end{itemize}

\begin{tcolorbox}[breakable,colback=teal!4,colframe=teal!35!black,title={Before you move on}]
What does \hi{ढलना} mean? (\textquotedblleft{}To set (of the sun), to wane\textquotedblright{}.) Say it once more.
\end{tcolorbox}
\section[\texorpdfstring{\hi{अधूरा} (adhūrā)}{अधूरा (adhūrā)}]{\hi{अधूरा} (adhūrā) — unfinished, incomplete}

\label{lesson:HI-C385-adhura}



\begin{quote}
Before the new one: say the Hindi for to tell off, then the Hindi for to set.
\end{quote}

\subsection*{You'll want to know: \hi{अधूरा}}
\textbf{\hi{अधूरा}} — \emph{adhūrā} — \textquotedblleft{}unfinished, incomplete\textquotedblright{}.

\begin{quote}
\hi{अगर} \hi{तुम} \hi{आज} \hi{नहीं} \hi{रुकोगे} \hi{तो} \hi{काम} \hi{अधूरा} \hi{रह} \hi{जाएगा।}
\end{quote}

\begin{quote}
\emph{agar tum āj nahīṃ rukoge to kām adhūrā rah jāegā.}
\end{quote}

\begin{quote}
If you don't stay today, the work will be left unfinished.
\end{quote}

\begin{cousinweb}
\textbf{\hi{अधूरा}} (\emph{adhūrā}) is inherited from \emph{*ardhapūraka}, \textquotedblleft{}half-full\textquotedblright{}, a form CDIAL reconstructs from its descendants (CDIAL 664).
\end{cousinweb}

\begin{grammarlens}[title={What you've built}]
One more word for saying what happens if.
\end{grammarlens}

\subsection*{Your turn}
\begin{itemize}
\raggedright
  \item \emph{Say it:} \emph{adhūrā}
  \item \emph{Say it:} \emph{adhūrā}, once more
  \item \emph{Say it:} \emph{agar tum āj nahīṃ rukoge to kām adhūrā rah jāegā.} — If you don't stay today, the work will be left unfinished.
  \item \emph{Recall:} say the Hindi for to tell off, then the Hindi for to set, then say \emph{adhūrā} again
\end{itemize}

\begin{tcolorbox}[breakable,colback=teal!4,colframe=teal!35!black,title={Before you move on}]
What does \hi{अधूरा} mean? (\textquotedblleft{}Unfinished, incomplete\textquotedblright{}.) Say it once more.
\end{tcolorbox}
